\chapter{Algoritmos implementados}
\label{cap:alg}

Para poder compararlos, todos los algoritmos siguen las mismas reglas
(\nolinkurl{src/algoritmos/__init__.py}). Reciben el grafo, un número máximo de evaluaciones
de la función de costo, una semilla aleatoria y un límite de tiempo. Devuelven la mejor ruta que
encontraron (el arreglo de vértices), su costo, cuántas evaluaciones usaron y el historial de cómo
fue mejorando. La semilla hace que cada corrida se pueda repetir y dé exactamente el mismo
resultado.

La tabla~\ref{tab:parametros} muestra los parámetros que usamos. Son valores iniciales tomados de
la investigación previa y de pruebas con grafos pequeños; no hicimos una calibración sistemática
sobre el FHCP (lo comentamos en las limitaciones).

\begin{table}[H]
\centering\small\filas
\caption{Parámetros usados en los experimentos. \emph{auto} quiere decir que el valor se calcula
según el tamaño del grafo o el presupuesto.}
\label{tab:parametros}
% Generado por Reporte/actualizar.py; no editar a mano.
\begin{tabularx}{\textwidth}{llL}
\toprule
Algoritmo & Archivo & Parámetros usados \\
\midrule
GA & \texttt{ga.py} & tam\_poblacion $=$ 40, prob\_cruce $=$ 0.9, prob\_mutacion $=$ 0.3, tam\_torneo $=$ 3, elitismo $=$ 2, prob\_inversion $=$ 0.5 \\
SA (swap) & \texttt{sa.py} & T0 $=$ auto, alfa $=$ auto, T\_min $=$ 1e-06, vecino $=$ swap \\
SA (mixto) & \texttt{sa.py} & T0 $=$ auto, alfa $=$ auto, T\_min $=$ 1e-06, vecino $=$ mixto \\
ACO & \texttt{aco.py} & hormigas $=$ 10, alfa $=$ 1.0, beta $=$ 2.0, rho $=$ 0.1, tau\_max $=$ 1.0, tau\_min $=$ auto, rotaciones $=$ 0, reinicio $=$ 100 \\
ACO + Pósa & \texttt{aco.py} & hormigas $=$ 10, alfa $=$ 1.0, beta $=$ 2.0, rho $=$ 0.1, tau\_max $=$ 1.0, tau\_min $=$ auto, rotaciones $=$ 50, reinicio $=$ 100 \\
ABC & \texttt{abc.py} & fuentes $=$ 20, limite $=$ auto, operador $=$ guiado \\
DE & \texttt{de.py} & tam\_poblacion $=$ 30, F $=$ 0.5, CR $=$ 0.9 \\
PSO & \texttt{pso.py} & particulas $=$ 30, w $=$ 0.729, c1 $=$ 1.49445, c2 $=$ 1.49445, v\_max $=$ 0.2, reinicio $=$ 30 \\
\bottomrule
\end{tabularx}

\end{table}

\section{Algoritmo genético (GA)}
\label{sec:ga}
Es nuestro algoritmo de control. Imita la evolución natural \cite{abc4}: hay una \textbf{población} de 40 rutas y en
cada generación se crean rutas nuevas combinando las mejores.
\begin{enumerate}
  \item \textbf{Selección por torneo:} se toman 3 rutas al azar y se queda la de menor costo.
  \item \textbf{Cruce OX} (\emph{order crossover}, con probabilidad 0.9): el hijo copia un tramo
        de un padre y completa el resto con los vértices que faltan, en el orden en que aparecen
        en el otro padre. Así el hijo siempre es una permutación válida.
  \item \textbf{Mutación} (probabilidad 0.3): intercambiar dos vértices o invertir un tramo.
  \item \textbf{Elitismo:} las 2 mejores rutas pasan directo a la siguiente generación, para no
        perderlas.
\end{enumerate}
El vértice 0 siempre va al inicio, porque una ruta y su rotación representan el mismo ciclo y así
evitamos tener soluciones repetidas en la población.

\section{Recocido simulado (SA)}
\label{sec:sa}
Se inspira en cómo se enfría un metal \cite{sa4}. Parte de una ruta al azar y en cada paso propone
un pequeño cambio:
\begin{itemize}
  \item Si el cambio mejora o empata el costo, se acepta.
  \item Si lo empeora, se acepta con probabilidad $e^{-\Delta/T}$, donde $\Delta$ es cuánto
        empeora y $T$ es la ``temperatura'' \cite{sa8}.
\end{itemize}
Al principio la temperatura es alta y se aceptan bastantes cambios malos, lo que ayuda a no
quedarse atorado. Después, la temperatura baja poco a poco ($T \leftarrow \alpha T$) y el
algoritmo se vuelve más exigente. Escogimos la temperatura inicial para que un empeoramiento de
una arista se acepte con probabilidad 0.8, y calculamos $\alpha$ para que la temperatura llegue a
su mínimo justo cuando se acaba el presupuesto. Así el enfriamiento es lo más lento posible, como
pide el enunciado.

Probamos dos versiones:
\begin{itemize}
  \item \alg{SA (swap)}: la versión literal del enunciado, que solo intercambia dos vértices.
  \item \alg{SA (mixto)}: la mitad de las veces intercambia y la otra mitad invierte un tramo de
        la ruta (2-opt). Invertir un tramo cambia solo dos conexiones, mientras que un
        intercambio cambia hasta cuatro, así que es un movimiento más ``suave''.
\end{itemize}

\section{Colonia de hormigas (ACO)}
\label{sec:aco}
Las hormigas reales dejan feromona en los caminos buenos y las demás tienden a seguirlos
\cite{aco1,aco3}. En nuestra versión, en cada iteración 10 hormigas construyen una ruta:
\begin{enumerate}
  \item Cada hormiga empieza en un vértice al azar y solo avanza por aristas que existen.
  \item Elige el siguiente vértice con una probabilidad que depende de dos cosas: la
        \textbf{feromona} de la arista y una \textbf{heurística} que prefiere los vértices con
        menos vecinos libres (regla de Warnsdorff). La idea es visitar primero los vértices
        ``difíciles'', antes de que se queden sin salida.
  \item Al final de la iteración, la feromona se evapora un 10\,\% y se refuerza en las aristas
        de la mejor ruta. Además la feromona tiene un mínimo y un máximo (variante MAX-MIN), y se
        reinicia si la búsqueda lleva 100 iteraciones sin mejorar.
\end{enumerate}

El problema es que en estos grafos las hormigas se atoran muy seguido: llegan a un vértice cuyos
vecinos ya fueron visitados. Por eso hicimos la variante \alg{ACO + Pósa}. Cuando una hormiga se
atora, en lugar de rendirse aplica una \textbf{rotación de Pósa} \cite{posa1976}: si el extremo
del camino tiene un vecino que está en medio del camino, se invierte la parte final para que el
extremo cambie y la hormiga pueda seguir avanzando. Se permiten hasta 50 rotaciones por hormiga.

Hay un detalle importante: estas rotaciones no las contamos como evaluaciones de la función de
costo, porque no calculan el costo de una ruta completa. Por eso, en los resultados, el número de
evaluaciones de \alg{ACO + Pósa} sale mucho más bajo que el de los demás, y no se debe comparar
directamente.

\section{Colonia artificial de abejas (ABC)}
\label{sec:abc}
Se basa en cómo las abejas buscan comida \cite{abc7,abc8}. Cada ``fuente de alimento'' es una ruta, y
usamos 20. En cada ciclo hay tres tipos de abejas:
\begin{enumerate}
  \item \textbf{Empleadas:} cada una intenta mejorar su fuente con un pequeño cambio. Si la ruta
        nueva es igual o mejor, la cambia.
  \item \textbf{Observadoras:} eligen fuentes al azar, pero con más probabilidad las de menor
        costo, y también intentan mejorarlas.
  \item \textbf{Exploradoras:} si una fuente lleva demasiados intentos sin mejorar ($2n$), se
        abandona y se reemplaza por una ruta nueva al azar. La mejor fuente nunca se abandona.
\end{enumerate}
El cambio que hacen las abejas es la \textbf{inversión guiada} de la sección~\ref{cap:repr}, que
repara conexiones rotas usando aristas reales del grafo. Otra decisión importante fue aceptar
los cambios que \emph{empatan} el costo: en estos grafos casi todos los movimientos dejan el
costo igual, y si solo aceptáramos mejoras estrictas, la búsqueda se quedaría congelada.

\section{Evolución diferencial (DE)}
\label{sec:de}
La evolución diferencial \cite{de10} fue creada para problemas con números reales, así que tuvimos que
adaptarla. Usamos \emph{random keys} \cite{de13} (sección~\ref{cap:repr}): cada solución es un vector de
números y la ruta sale de ordenarlos. Con una población de 30 vectores, en cada generación y para
cada vector $x$:
\begin{enumerate}
  \item \textbf{Mutación:} se crea un vector nuevo $v = x_{r_1} + F\,(x_{r_2}-x_{r_3})$ a partir de
        tres vectores distintos elegidos al azar, con $F=0.5$.
  \item \textbf{Cruce:} cada componente se toma de $v$ con probabilidad $CR=0.9$, y si no, de $x$.
  \item \textbf{Selección:} si la ruta del vector nuevo cuesta lo mismo o menos que la de $x$, lo
        reemplaza.
\end{enumerate}
En la investigación encontramos otras formas de adaptar DE a permutaciones, como la regla SPV
\cite{de17}. Escogimos \emph{random keys} porque es la más sencilla y garantiza que nunca salga
una ruta inválida.

\section{Enjambre de partículas (PSO)}
\label{sec:pso}
Simula un grupo de partículas que se mueven buscando el mejor lugar; cada una recuerda su mejor
posición y conoce la mejor de todo el enjambre \cite{pso1995}. Igual que en DE, cada partícula es
un vector de \emph{random keys}. Usamos 30 partículas, y en cada paso la velocidad $v$ y la
posición $x$ de cada una se actualizan así:
\[
v \leftarrow w\,v + c_1 r_1 (p - x) + c_2 r_2 (g - x), \qquad x \leftarrow x + v,
\]
donde $p$ es la mejor posición de la partícula, $g$ la mejor del enjambre y $r_1,r_2$ son números
al azar entre 0 y 1. El término $w$ es el \emph{peso de inercia} \cite{pso1998}, que controla
cuánto conserva la partícula su dirección anterior. La velocidad se limita a $\pm0.2$ para que
las partículas no salgan disparadas.

Al probarlo vimos que el enjambre se juntaba muy rápido alrededor de la mejor partícula y dejaba
de explorar, incluso con grafos de 8 vértices. Por eso agregamos un \textbf{reinicio}: si la mejor
solución no mejora en 30 pasos, todas las partículas, excepto la mejor, se mandan a posiciones
nuevas al azar.
